\documentclass[11pt]{article}
\usepackage[margin=0.85in]{geometry}
\usepackage{amsmath,amssymb,booktabs,longtable,xurl,hyperref}
\usepackage[T1]{fontenc}
\hypersetup{colorlinks=true,urlcolor=blue}
\setlength{\emergencystretch}{2em}
\title{Restore CLF Priority and Post-Avoidance Nominal Recovery}
\author{Pan Dynamics Collision Avoidance Research}
\date{October 1, 2026}
\begin{document}
\maketitle

\section{Problem and implementation}
The preceding causal study, \path{CLF_CAUSAL_DIAGNOSIS_20260930.tex},
found complete zero-CLF-slack trajectories at states where the controller
instead issued an energy-increasing input. Whole-trajectory backtracking
removed useful first-input descent to repair the nonlinear terminal core.
Any small same-state slack improvement then counted as completed CLF work.
A weighted squared-slack objective could also prefer positive relaxation
over available zero relaxation. Finally, shifting consumed the initial
completion allowance until a short straight-terminal countdown reached the
issued input, even during unobstructed circular cruise.

The implementation changes the numerical realization of the same predictive
controller in \path{controller/solvePredictiveControl.m}. The target model,
continuous vehicle dynamics, nominal path, collision buffer and physical
admission constraints are retained. The existing RK4/safety/terminal mesh
is refined consistently as described below. Road boundaries remain absent.

\begin{enumerate}
\item \textbf{Give actual CLF slack priority.} An admitted zero-safety
witness establishes the primary PCBF minimum. If its CLF slack is positive,
optimization continues. A conic subproblem first seeks zero CLF slack, with
tracking and input-deviation costs inside that slice. If the slice cannot
be solved, the relaxed subproblem minimizes CLF slack alone. Positive slack leaves later inputs undetermined; a subsequent
predictive-cost solve fixes the minimizing first input and optimizes the
later trajectory without changing that first-step nonlinear CLF value. The former
\texttt{clf.relaxationWeight} is removed and rejected by configuration
validation. Input proximity is measured from the current linearization
input, never from zero.
\item \textbf{Preserve first-input descent during nonlinear correction.}
Analytic first-step Gauss--Newton directions seek the actual CLF sublevel.
Their input displacement is bounded by the existing local input trust box;
relinearization, rather than a single extrapolated Newton jump, supplies
larger changes. Unbounded correction steps in an intermediate experiment
reached extreme combined-slip inputs and produced an ill-conditioned
prediction during a later encounter.
Later-input correction fixes the first control while repairing terminal
membership. When other constraints also need correction, a condensed
minimum-input-deviation QP includes sampled collision, physical, slew and
terminal-separation rows. Every trial is replayed through the nonlinear
model and passes the original hard/zero-safety-slack admission before it
can issue a control. Raw and corrected candidates remain independently
eligible. These operations are within the existing optimizer.
\item \textbf{Repair nonlinear rollout defects.} A small measured-successor
discrepancy can grow along an unstable open-loop tail until its tire model
is no longer evaluable. The same issue can invalidate a conic proposal
before endpoint correction becomes available. A defined rollout can also
accumulate enough error to destroy terminal or collision feasibility.
An unevaluable proposal, or an inadmissible full conic proposal with an
admissible first stage, is projected
onto nonlinear dynamics using finite-horizon linear-quadratic feedback
around its proposed state path. The first input is fixed; only later
inputs compensate rollout error. Gains are built along the evaluable
anchor using existing numerical state scales and input weights.
The resulting trajectory is a search candidate and must pass the same
complete nonlinear admission.
\item \textbf{Use a meaningful stopping condition.} The global
nonnegative CLF lower bound is attained only by an admitted trajectory
whose actual residual is zero, apart from floating-point roundoff.
Positive local stationarity is reported separately. A time or iteration
limit retains the best admitted witness and reports incomplete CLF work.
A mere improvement over the incumbent is no longer sufficient.
\item \textbf{Maintain the moving completion allowance.} The prediction
length may shorten to
\[
H_{\min}=\min\{H_{\max},N+\lceil T_{\rm recovery}/h\rceil\},
\]
then grows its endpoint with the already admitted terminal policy as the
window advances. The existing three-second allowance is preserved. It is
a prediction length, not a required recovery time or fixed endpoint pose.
The shift-and-append recursive-feasibility construction is unchanged.
\end{enumerate}

The zero-slack cone is represented directly as
\(\|L\widehat e_1\|_2\le\sqrt{(1-\alpha)W(e_0)}\), avoiding cancellation
between nearly equal rotated-cone constants near trim. Positive scaling
of the CLF and tracking expressions improves numerical conditioning
without changing feasible sets or objective ordering. Endpoint corrections
cannot modify the actual first successor they are intended to preserve.
The added correction QP has at most 100 active-set iterations; up to five
nonlinear correction steps share the existing soft call budget.
For a conic step of size $d$ near a regular anchor, its affine dynamic
defect is second order. The projected rollout error has the local form
\[
\epsilon_{j+1}=(A_j+B_jK_j)\epsilon_j+O(\|d\|^2),
\qquad \widetilde u_0=u_0.
\]
This is a numerical retraction that preserves the proposed first-step
CLF value. It does not impose a scheduled path position, change the
terminal set or establish a new closed-loop robustness theorem.

\section{Conditional recovery statement}
The nominal error contains lateral position, heading relative to curved
trim, body longitudinal/lateral velocity, and yaw rate. It does not contain
a scheduled longitudinal position. The positive-definite LQR metric is
\(W(e)=e^TPe=\|Le\|_2^2\). The controller evaluates
\[
\rho_k=\max\{0,W(e_{1|k})-(1-\alpha)W(e_k)\},\qquad\alpha=0.01.
\]
For the declared nominal transition, if zero relaxation continues to be
achieved after some index $k_0$, then
\[
W(e_k)\le(1-\alpha)^{k-k_0}W(e_{k_0})\longrightarrow0.
\]
Positive definiteness implies convergence of all five nominal errors,
hence given-path constant-speed cruising. No settling-time requirement is
introduced. The original local Riccati construction is retained; this
work does not prove that its metric is a global nonlinear CLF.

The implementation reports the raw residual and uses only
\(64\operatorname{eps}(\max\{1,W_k,W_{k+1}\})\) for lower-bound stopping.
This entails a tiny practical numerical bound rather than an exact-real
arithmetic claim. A positive local stationary result does not prove that
zero relaxation is globally infeasible. A finite-budget local solver
cannot guarantee finding every disconnected feasible zero-slack solution.
Absence of an obstacle alone is not substituted for a successful actual
CLF/trajectory check. Likewise, the RK4 statement is distinct from the
tightly integrated ODE45 experimental plant.

\section{Validation protocol and results}
The final production source hashes are recorded with the compact evidence
in \path{CLF_NOMINAL_RECOVERY_20261001/}. The fourteen original threat
fixtures, at 8 and 15 m/s, are unchanged. Every fixture starts separated
and collides if the ego follows its given path at constant speed.
The hold is 50 ms; prediction uses the existing RK4 map and a 6-mm
collision buffer. Tight ODE45 supplies measured ego successors with
relative/absolute tolerances $10^{-11}$/$10^{-12}$; independent rectangle
geometry checks 31 times per hold. There is no observer, noise, delay,
random draw or road constraint.

An initial implementation-only campaign exposed short corner contacts
between the former 8.33-ms collision-check nodes. For example, the
15-m/s turning-crossing replay overlaps at 1.860--1.865 s while both
neighboring prediction nodes satisfy the buffer. The final default
integration step is 5 ms, giving ten intervals per hold, consistently
refining the dynamics, finite collision checks and terminal enclosures.
A second campaign showed that mesh refinement alone could still miss a
short corner contact. The final admission therefore also bounds the
rectangle distance on every interval of the nominal pose interpolant.

Let $d_a,d_b$ be the endpoint distances on an interval of duration $\tau$,
and let $L_d$ bound the relative translational speed plus the two angular
speeds multiplied by their respective body radii. Distance is Lipschitz,
so a valid lower bound is
\[
\underline d=\min\{d_a,d_b,(d_a+d_b-L_d\tau)/2\}.
\]
An interval is accepted when this bound exceeds half the declared sampled
margin. Otherwise, it is bisected and the new midpoint must satisfy the
\emph{full} sampled margin. At the default 6-mm sampled margin the
admitted nominal interpolant thus has a strictly positive 3-mm interval
bound. The relative-speed construction cancels common translation.
Additional midpoint rows are supplied to the same optimization and
nonlinear correction problems. Curved terminal enclosures include the
nominal interpolation chord error.

An intermediate implementation also exposed an unchecked-prefix cache
defect: a not-yet-validated interval could survive endpoint correction as
if it had been checked. All target-bearing intervals now start as
uncertified, and only actually checked unchanged prefixes may be reused.
The measured-successor regression independently reconstructs every
returned plan to guard this admission property. A shifted trajectory
that loses admission is first offered to the same later-input correction
before a full conic problem is rebuilt.

No audit threshold or physical collision buffer was weakened. This
interval argument concerns linear interpolation of the nominal RK4 ego
poses and the exact prescribed target flow. It is not an error bound
between that interpolant and the physical continuous-time plant.

Each threat run stops after all five absolute errors remain below
\[
[0.10\,\mathrm m,1^\circ,0.10\,\mathrm{m/s},0.05\,\mathrm{m/s},
0.01\,\mathrm{rad/s}]
\]
for five seconds, no earlier than eight seconds.
The observation cap is 80 seconds, not a required settling time. A finite
dwell is an observed recovery criterion, not an asymptotic proof.
Two additional no-target controls retain exact circular cruise for 60
seconds. Behavior tests also start with a 0.1-m lateral perturbation and
check geometric CLF decay beyond the original initialization-tail length.

\begin{longtable}{llrrr}

\toprule

Speed (m/s) & Scenario & Prior (s) & Current (s) & Gap (mm) \\

\midrule

\endhead

8 & acceleratingHeadOn & 725.90 & 13.05 & 66.286 \\

8 & brakingLead & 793.10 & 14.80 & 6.002 \\

8 & crossing & 559.75 & 13.05 & 6.942 \\

8 & curvedCrossing & >300 & 17.90 & 9.867 \\

8 & curvedHeadOn & >300 & 19.60 & 26.029 \\

8 & headOn & 670.05 & 13.05 & 142.060 \\

8 & turningCrossing & 556.45 & 16.00 & 24.696 \\

15 & acceleratingHeadOn & 19.55 & 14.75 & 11.312 \\

15 & brakingLead & 36.10 & 16.15 & 5.976 \\

15 & crossing & 47.35 & 13.05 & 6.630 \\

15 & curvedCrossing & >300 & 16.45 & 7.479 \\

15 & curvedHeadOn & >300 & 13.05 & 6.497 \\

15 & headOn & 20.80 & 13.05 & 88.310 \\

15 & turningCrossing & 60.05 & 15.60 & 12.344 \\

\bottomrule

\end{longtable}

All fourteen final closed-loop runs complete without a controller failure, preserve positive independently sampled clearance, and satisfy the observed nominal-recovery dwell. The full 6-mm replay buffer is satisfied in 13 of fourteen runs.

The below-buffer results are 15 m/s brakingLead (5.97552 mm). They remain collision-free; a 6-mm prediction-sample constraint is not claimed as an identical ODE45 clearance bound.

Maximum final lateral error over the fourteen stopping states is 0.032547 m. The reported zero-slack flag agrees with the actual roundoff-only CLF predicate in every trace.

The largest controller-call time is 5.539001 s (8 m/s turningCrossing). Across cases, subsequent-call medians range from 0.036740 to 0.127972 s. These times include formulation, nonlinear checking and correction, not only conic optimization.

The maximum absolute lateral error anywhere in a final threat run is 4.445576 m. Final zero-CLF-slack suffixes last at least 5.05 s. The largest positive independently observed ODE successor excess over the predicted 0.99-W bound, normalized by max(1,W), is 1.49127e-07. This difference is retained rather than absorbed into the reported CLF slack.

Across the fourteen runs, 536 holds retain an admitted safety witness while reporting incomplete CLF optimization, and 8 holds report positive local stationarity. These are not reported as attained zero-slack optima or globally infeasible zero-slack problems.

\begin{center}\begin{tabular}{rrrr}\toprule Speed & No-target duration (s) & Maximum $W$ & Final $|e_y|$ (m) \\ \midrule

15 & 60.0 & 8.75e-20 & 2.81e-12 \\

8 & 60.0 & 6.96e-21 & 8.75e-12 \\

\bottomrule\end{tabular}\end{center}

All 305 targeted MATLAB tests in thirteen files and all 15 independent Python audit tests pass. Factory Code Analyzer findings are retained in the bundle; performance advisories remain. The unrelated full repository suite was not run.


\section{Scope, residual limitations and reproduction}
The collision audit separately records strict positive clearance and the
configured 6-mm buffer. A prediction-sample buffer does not establish the
same minimum at every independent ODE45 sample or at every continuous
time. All measured discrepancies are retained; the audit threshold is
not weakened to make a report pass.

The opposing circular-target and accelerating turning-crossing fixtures
can have recurring conflicts by their original definitions. In particular,
the 8-m/s turning-crossing target approaches again around nine seconds;
it has not permanently departed after the initial crossing.
Recovering nominal cruise for the observed
dwell does not establish permanent nominal cruise through all later
opposing encounters. The algorithm retains the original fixed target
epoch and indefinite terminal separation test, without assuming one-time
departure or editing those fixtures.

Concurrent single-thread MATLAB workers provide functional validation,
not an isolated timing benchmark. The existing five-second budget remains
soft: an in-flight factorization or validation may finish beyond it.
Some conflict frames still exhaust that budget. This work establishes
the reported recovery improvement, not a 50-ms real-time guarantee.

The compact bundle includes per-case and per-second metrics, independent
audits, final test results, source/raw hashes and actual reproduction
commands. Full JSON/MAT trajectories and generated PDFs remain in the
external research cache. Preliminary prototypes and interrupted runs are
identified separately and are not counted as final validation. Unrelated
observer work and external solver dependencies are excluded from this
change.
\end{document}
